% titel: Stosse: Elastische und inelastische Stösse
% skript: Kap. Stösse, Impuls- und Energieerhaltung

In dieser Aufgabe werden zentrale und nicht-zentrale elastische Stösse sowie
die inelastische Verformungsenergie bei Verkehrsunfällen betrachtet. Gemeinsames
Werkzeug sind Impuls- und (je nach Stossart unterschiedlich) Energieerhaltung.

\begin{werkzeugkasten}[Werkzeuge]
\textbf{Zentraler elastischer Stoss} (Massen $m_1$ (mit $v_1$), $m_2$ (ruhend)):
Geschwindigkeiten nach dem Stoss
\begin{equation}
  v_1' = \frac{m_1-m_2}{m_1+m_2}\,v_1\;,\qquad
  v_2' = \frac{2m_1}{m_1+m_2}\,v_1\;.
\end{equation}

\textbf{Übertragene Energie} auf die ruhende Masse $m_2$:
\begin{equation}
  \frac{E_2'}{E_1}
  = \frac{4\,m_1 m_2}{(m_1+m_2)^2}\;.
\end{equation}

\textbf{Mehrfache Stösse:} Nach $N$ identischen Stoessen eines Teilchens der
Masse $m_1$ gegen ruhende $m_2$ bleibt der Bruchteil
\begin{equation}
  \left(\frac{m_1-m_2}{m_1+m_2}\right)^{2N}
\end{equation}
der Anfangsenergie beim Mittelpunkt-Teilchen.

\textbf{Nicht-zentraler elastischer Stoss gleicher Massen:} Die beiden
Teilchen fliegen nach dem Stoss im rechten Winkel zueinander davon; die
Impulskomponente senkrecht zur Verbindungslinie bleibt beim jeweiligen
Teilchen erhalten.

\textbf{Inelastischer Stoss (vollplastisch):} Impuls erhalten, kinetische
Energie nicht. $\Delta E$ steht für Deformation zur Verfügung; sie ist der
Verlust an kinetischer Energie.
\end{werkzeugkasten}

\begin{werkzeugkasten}[Vorgehen in allgemeiner Form]
Elastische zentrale Fälle: $v_2'$ bzw. Energietransfer aus den beiden
Standardformeln. Vielfach-Stoss: Potenzgesetz anwenden. Nicht-zentraler Fall:
Kinematik im rechten Winkel. Unfälle: $\Delta E_\mathrm{Verformung}$ aus
Impuls- und Energiebilanz.
\end{werkzeugkasten}

\begin{achtung}[Typische Fallen]
\begin{itemize}
  \item Beim Translations-Impuls gilt die Erhaltung des \emph{Impulsvektors},
       nicht des Betrags; Richtungen separat ansetzen.
  \item Beim nicht-zentralen Stoss ist die \emph{Gesamtenergie} (nicht jede
       Komponente) erhalten; die Rechnung über den rechten Winkel vereinfacht.
  \item Bei Verkehrsunfällen $v$ in \unit{\kilo\metre\per\hour} erst in
       \unit{\metre\per\second} umrechnen.
\end{itemize}
\end{achtung}
